\subsubsection{Min25}
Min25 示例：积性函数满足 $f(p^e)=p^e(p^e-1)$，求 $\sum_{i=1}^n f(i)$ 模 $10^9+7$。一般函数需要替换质数处点值和及素数幂处取值，不能只把一般积性函数拆成完全积性函数后直接组合答案。定义 $S(m,j)$ 为不含前 j 个素因子的 $2..m$ 的函数和，最终答案是 $S(n,0)+1$。数组按 $\sqrt n$ 和实际商块数动态分配，修复原固定容量无法覆盖所声称 $10^{11}$ 范围的问题。每次 sum(n) 重新初始化；时间为 Min25 筛及素数幂递归的实际开销，不承诺固定秒数。
\begin{lstlisting}
struct Min25 {
    static constexpr ll mod=1000000007;
    ll n;int lim;
    vector<int> primes,small,large;
    vector<ll> values,g1,g2,s1,s2;
    int index(ll x)const{return x<=lim?small[x]:large[n/x];}
    ll power_value(ll x)const{x%=mod;return x*((x-1+mod)%mod)%mod;}
    ll S(ll m,int j){ // exclude the first j primes
        if(m<2||(j&&primes[j-1]>=m))return 0;
        int id=index(m);
        ll ans=(g2[id]-g1[id]-s2[j]+s1[j])%mod;
        if(ans<0)ans+=mod;
        for(int i=j;i<int(primes.size())&&ll(primes[i])*primes[i]<=m;i++){
            ll p=primes[i],v=p;int exponent=1;
            for(;;){
                ans=(ans+power_value(v)*(S(m/v,i+1)+(exponent>1)))%mod;
                if(v>m/p)break;v*=p;exponent++;
            }
        }
        return ans;
    }
    ll sum(ll input){
        assert(input>=0&&input<=100000000000LL);
        if(input<1)return 0;n=input;lim=sqrtl(n);
        while(ll(lim+1)<=n/(lim+1))lim++;
        while(ll(lim)>n/lim)lim--;
        primes.clear();vector<char> composite(lim+1);
        for(int i=2;i<=lim;i++){
            if(!composite[i])primes.push_back(i);
            for(int p:primes){if(ll(i)*p>lim)break;composite[i*p]=true;if(i%p==0)break;}
        }
        s1.assign(primes.size()+1,0);s2=s1;
        for(int i=0;i<int(primes.size());i++){
            ll p=primes[i];s1[i+1]=(s1[i]+p)%mod;s2[i+1]=(s2[i]+p*p)%mod;
        }
        small.assign(lim+1,0);large=small;values.clear();g1.clear();g2.clear();
        for(ll l=1,r;l<=n;l=r+1){
            ll v=n/l;r=n/v;int id=values.size();values.push_back(v);
            if(v<=lim)small[v]=id;else large[n/v]=id;
            ll x=v%mod;
            g1.push_back((x*(x+1)%mod*500000004%mod+mod-1)%mod);
            g2.push_back((x*(x+1)%mod*((2*x+1)%mod)%mod*166666668%mod+mod-1)%mod);
            if(r==n)break;
        }
        for(int i=0;i<int(primes.size());i++){
            ll p=primes[i];
            for(int j=0;j<int(values.size())&&p*p<=values[j];j++){
                int k=index(values[j]/p);
                g1[j]=(g1[j]-p*((g1[k]-s1[i]+mod)%mod)%mod+mod)%mod;
                g2[j]=(g2[j]-p*p%mod*((g2[k]-s2[i]+mod)%mod)%mod+mod)%mod;
            }
        }
        return (S(n,0)+1)%mod;
    }
};
\end{lstlisting}

\paragraph{使用条件与多测} 本实现范围 $0\le n\le10^{11}$，仍须按实际内存和时限判断是否可行，$O(\sqrt n)$ 存储；不是针对完整 64 位区间设计的接口。素数前缀表可针对固定最大上界复用，但本实现选择逐次重建以明确多测语义。
